\chapter{The components: V°Edit, V°Text, V°Protect}\label{ch:ghcomponents}

miniVivisection's Grasshopper part, \file{miniVivisection.GH.gha}, adds three components on the tab
\ui{miniVivisection} -- \ui{V°} when Grasshopper's tabs don't fit: Grasshopper shortens tabs from the right
until they do (with the mark chosen for the V names, \see{sec:vnames}). At \menu{miniVivisection\>Help} it
adds V°Examples, V°Manual and V°Settings too (\see{sec:findexamples}).

\section{V°Edit: Grasshopper geometry edited in Rhino}\label{sec:vedit}

Tab \menu{miniVivisection\>Debug}. Some geometry is easier to fix by hand than with numbers. V°Edit sits in a
definition between two components, puts chosen items of its input into Rhino as ordinary objects, and feeds your
edits back into the definition. They are saved with the \file{.gh}.

\shot{11_v_edit}{Lesson GH~4: V°Edit between \emph{Profiles} and \emph{Loft} puts profile [1] into Rhino as an edit object
  (\emph{1 in Rhino}), with the name dot \emph{V°Edit [1]}. \ui{Reset sel.} and \ui{Reset all} drop edits.}{fig:vedit}

\subsection{Inputs and outputs}

\begin{description}
  \item[Geometry] an item, a list or a tree.
  \item[Index] which items to edit, any mix of:
    \begin{itemize}
      \item \textbf{Indices}: positions in every branch; negative counts from the end. Or per branch, with an index
            \emph{tree} matched by path;
      \item \textbf{Paths}: \code{\{0;2\}} as text, or a Path parameter: only those branches, all their items, or with
            indices the items at them;
      \item \textbf{Path with index}: \code{\{0;2\}[3]}, that one item.
    \end{itemize}
    Nothing connected: item 0.
  \item[Outputs,] in the input's exact tree: \emph{Geometry}, with your edits in place, and \emph{Edited}, True or False
        for every single item.
\end{description}

\subsection{Editing}

\begin{enumerate}[label=\arabic*.]
  \item The items at \emph{Index} show in Rhino as ordinary objects, \textbf{blue}, each with a name dot and tail showing
        the component's nickname. Rename the component, and the dots follow at once.
  \item Edit them with anything: Gumball, control points, Move, Scale, Rebuild.
  \item When you finish, the edited geometry \textbf{replaces} its item, and the definition recomputes from the component on.
  \item The object turns \textbf{red-orange}, its dot says \emph{edited}, the component counts the edits, and a balloon
        on it explains.
\end{enumerate}

Every later solution uses your edits instead of the incoming items, also when the input changes. When the input
changes so that another kind of item is at an edited place, the edit isn't used there: a warning names it, and
\ui{Reset all} drops it.

\subsection{Undo, reset, save}

\begin{itemize}
  \item \ui{Reset sel.} drops the edits of the edit objects selected in Rhino (orange while one is selected);
        \ui{Reset all} drops every edit (orange while there are any). The menu's \ui{Reset selected edits} and
        \ui{Reset all edits (n)} do the same.
  \item Undo in Rhino back to the input's geometry drops an edit; Grasshopper's Undo takes back an edit made in
        Rhino. Grasshopper's save prompt includes the edits.
  \item The menu has \ui{Select edit objects in Rhino}, \ui{Edit objects in Rhino} (show or hide them), \ui{Name dots in Rhino},
        and \ui{About V°Edit…}, with its version, build and build time, and the plug-in's.
  \item \emph{Preview}: Grasshopper's own preview of the Geometry output, as for other components.
\end{itemize}

\begin{note}[The edit objects belong to the component]
They are not part of the model: they go away while Grasshopper is minimized or hidden, another definition is on
the canvas, or the component is disabled; when it is deleted or the definition closed. They are never saved into
the \file{.3dm}, hidden or locked ones neither, and don't mark the model as changed. (Export Selected exports them
when they are selected.)
\end{note}

\section{V°Text: a Panel's text as Rhino text}\label{sec:vtext}

Tab \menu{miniVivisection\>Text}. V°Text makes Rhino text of a Panel's text, the way the Panel shows it: Markdown laid
out, a script as code, plain text with its lines kept. Notes and scripts become drawing annotation in the model, in
the format they have on the canvas.

\shot{12_v_text}{A Markdown \emph{Notes} Panel and a Python \emph{Ring script} Panel, each feeding a V°Text. Rhino shows the
  Markdown laid out and the script as code. The red frames are the \emph{Bounds} outputs, the crosses the anchors.}{fig:vtext}

\subsection{What it makes}

\begin{description}
  \item[Markdown] as \code{MarkdownText} lays it out (\cref{ch:markdowntext}): headings, bold, italic, code, lists, task
        boxes, quotes, code blocks, tables, rules, links with their bars.
  \item[A Python or C\# script] as code: Consolas, line numbers in a gutter, a frame, each token in its syntax color (one
        text object per color, lined up).
  \item[Anything else] as plain text, its lines kept.
\end{description}

\ui{Show as} in its menu: \ui{Autodetect (code or Markdown)} follows the Panel the text comes from, its own
\ui{V° Show as}, else what it is; or \ui{Markdown (rendered)}, \ui{Python code}, \ui{C\# code}, \ui{Plain text}.

\subsection{Inputs}

\begin{xltabular}{\linewidth}{@{}P{0.16\linewidth}L@{}}
\toprule
\ui{Text}   & connect a Panel; a list is taken as its lines \\
\ui{Height} & a paragraph's or a line's text height (Rhino's text height: capital letters). \textbf{10~pt} by default: points made the model's units (10~pt = 3.53~mm), or model units (menu: \emph{Sizes in points / in model units}) \\
\ui{Anchor} & a point; the menu chooses the text's corner at it: \emph{top left} (the default), \emph{top right}, \emph{bottom left}, \emph{bottom right} \\
\ui{Angle}  & turned about the anchor in the world's XY plane (radians; \emph{Degrees} in the input's menu) \\
\ui{Width}  & the width Markdown and plain text wrap in, in the Height's unit; 0: no wrapping. Without a Width it wraps about as wide as the Panel (plain text only when the Panel wraps); making the Panel wider or narrower, or changing its Show as, updates it \\
\ui{Font}   & the base font of Markdown and plain text (code stays in Consolas); nothing: the dimension style's. Right-click \ui{Set one Text} lists the installed fonts, each shown in itself: type to filter, double-click or OK to take one \\
\bottomrule
\end{xltabular}

\subsection{Outputs}

\begin{description}
  \item[Text] the Rhino text objects as the component draws them, a token or a part each.
  \item[Decorators] everything else it is drawn with: the lines of code frames, rules, quotes' bars, tables' grids, and
        the links' bars (hatches). Both placed and turned as the text is; their previews are off, since the component draws
        the text in its colors.
  \item[Bounds] the text's bounding rectangle, aligned with it.
  \item[Anchor] the anchor point.
\end{description}

The component's message says what it shows and where it is anchored (\emph{Markdown, top left}). The preview draws the text
in Grasshopper's preview color, code in its syntax colors, green while selected.

\paragraph{Bake} (right-click \ui{Bake}) makes Rhino objects, grouped; a token's color is its object's color. Baked Markdown
keeps what \code{MarkdownText}'s parts keep: double-click any of it to edit its Markdown in place, and its links work. Links in
the preview aren't clickable: bake it.

\section{V°Protect: is it still the original?}\label{sec:vprotect}

Tab \menu{miniVivisection\>Files}. Lessons, templates and handed-out definitions get changed by accident. V°Protect tells
whether a definition, or other files, are still the originals, and puts them back. It has no inputs and no outputs; its
menu sets everything, saved with the definition.

\subsection{Setting it up}

\begin{description}
  \item[Files:] None set: the definition's own file (\ui{Files: this definition}). In its submenu: \ui{Add files...}
        (several at once), \ui{Add this definition}, each file's \ui{Remove}, \ui{Remove all (keep this definition)}.
  \item[Where the originals go:] \ui{Archive}: by default \file{(the definition's name).protect.zip} next to it;
        \ui{Choose...} another (an archive that exists is added to), or \ui{Next to the definition (the default)}. Or
        \ui{Keep the originals in the definition itself}: no archive. The other files' originals are kept in the component, saved with the definition, which grows by
        their size; the definition keeps its own original by a fingerprint of every object's saved state.
  \item[Store now:] Stores at once. A definition not saved yet can't be kept; it says so.
\end{description}

In setup, the files are hashed (MD5) and copies stored in the archive whenever the component solves: placed, opened, a
setting changed, \ui{Store now}. Two files of the same name in different folders are two entries; a file changed since is
replaced, an unchanged one left as it is. Its \textbf{message balloon} says per file \emph{added}, \emph{replaced},
\emph{unchanged} or \emph{missing} (a warning); its tooltip lists the files, the archive and the whole log.

\subsection{Protect and Revert}

\begin{description}
  \item[Protect] (the button under it, or its menu): the definition is saved and stored as the original, with the component
        protecting it. From then on the component looks like a Grasshopper Button: in front it says \textbf{original}, the
        component \textcolor{mvgreen}{green}, while the file on disk is the original and the open definition has no changes;
        \textbf{modified}, the component \textcolor{mvred}{red}, when either isn't so. Then a gray \textbf{Revert} button. Its
        balloon and tooltip say which, and how the other files stand.
  \item[Revert] (the button while \emph{modified}, or \ui{Revert to the original} in its menu) asks first. When the file
        on disk isn't the original, the archived copy is written over it; then the definition is opened again from the
        file, and the changes of the open one are dropped.
  \item[Back to setup] (its menu: right-click anywhere on it) ends the protection: its files and where their originals
        go can be changed again.
\end{description}

\paragraph{In a Panel's Markdown} \code{\{\{protected\}\}} shows \emph{protected} or \emph{not protected},
\code{\{\{modified\}\}} \emph{modified}, \emph{not modified} or \emph{not tracked} (not protected): the definition's
V°Protect's state, following it as it changes. A link \code{[Revert](protect:revert)} is its \ui{Revert}, asked first.
V°Text shows the fields the same way. For Rhino files: \code{ProtectModel} (\see{sec:protectmodel}).

\begin{note}
V°Protect has no password and no write protection: it tells, and it puts back.
\end{note}

\begin{lesson}{the overview and every Grasshopper lesson}
\lessonfile{gh0_overview.gh} has a V°Protect to try: move its slider, it turns red; \ui{Revert} puts the overview back. Every
Grasshopper lesson has one under its text, keeping the original in the lesson itself. \lessonfile{gh4_v_edit}: V°Edit.
\end{lesson}
